\FifthArticle{17}
The procedure that must be used to determine the imaginary roots is quite different from that taught in the preceding articles. In general, finding imaginary roots by an indirect method is far more difficult than finding real roots, because the former must be sought in an infinite region of two dimensions, the latter in an infinity of only one dimension. Precisely for this reason, a very broad special case in which that difficulty can be avoided and the question brought into the same domain as the search for real roots deserves a separate treatment. Equations with three terms provide such a case.

Since the method can be applied with equal ease whether the coefficients of the equation are real or imaginary, I take at once the general form of the equation as my starting point:
\[
\mathrm{X}=x^{m+n}+e(\cos\varepsilon+i\sin\varepsilon)x^m+f(\cos\varphi+i\sin\varphi)=0
\]
where $e$ and $f$ denote positive quantities. For a real coefficient, positive or negative, the corresponding angle ($\varepsilon$ or $\varphi$) is then either $0$ or $180^{0}$. Here too the assumption that $m$ and $n$ have no common divisor can be retained without loss of generality. An imaginary root $x=t+iu$ satisfying the equation is put in the form $r(\cos\rho+i\sin\rho)$. For our present purpose
\SourceLeaf{104}
it is more advantageous to omit the usual condition that $r$ be positive and require instead that $\rho$ always be taken between the limits $0$ and $180^{0}$. When both coefficients of the equation are real, the range of values of $\rho$ can be further reduced by half. For, as is known, the imaginary roots of such an equation belong together in pairs, such as $t+iu$ and $t-iu$. It is then evident that, for one root of each pair, the value of $\rho$ falls between $0$ and $90^{0}$, and only these need be determined by the indirect procedure: the other follows of itself by exchanging $\rho$ for $180^{0}-\rho$ and $r$ for $-r$.

\FifthArticle{18}
The essence of the method is to set up an equation containing only $\rho$, without $r$. To obtain it, divide the equation $\mathrm{X}=0$ by its first term, putting it in the form
\[
1+e(\cos\varepsilon+i\sin\varepsilon)x^{-n}+f(\cos\varphi+i\sin\varphi)x^{-m-n}=0
\]
or
\[
\begin{aligned}&1+er^{-n}(\cos(n\rho-\varepsilon)-i\sin(n\rho-\varepsilon))\\
&\quad+fr^{-m-n}(\cos[(m+n)\rho-\varphi]-i\sin[(m+n)\rho-\varphi])=0\end{aligned}
\]
Since the imaginary parts here must cancel, one has\nobreak\hspace{.5em}(I)
\[
r^m=-\frac{f\sin((m+n)\rho-\varphi)}{e\sin(n\rho-\varepsilon)}
\]
Similarly, dividing the equation $\mathrm{X}=0$ by its second or third term and observing that in both cases the imaginary parts of the new equations must cancel gives the equations
\[
\left.
\begin{array}{rcl}
r^{m+n}&=&\displaystyle\frac{f\sin(m\rho+\varepsilon-\varphi)}{\sin(n\rho-\varepsilon)}\\[1.2ex]
r^n&=&\displaystyle-\frac{e\sin(m\rho+\varepsilon-\varphi)}{\sin((m+n)\rho-\varphi)}
\end{array}
\right\}\quad\mathrm{(I)}
\]
Each of the three equations (I) can evidently also be derived by combining the other two. But eliminating $r$ from a combination of two gives (II)
\SourceLeaf{105}
\[
\lambda=(-1)^{m+n}\frac{\sin(m\rho+\varepsilon-\varphi)^m\sin(n\rho-\varepsilon)^n}{\sin((m+n)\rho-\varphi)^{m+n}}
\]
where, for brevity, as above,
\[
\frac{f^n}{e^{m+n}}=\lambda
\]
is put. The various values of $\rho$ must be determined from this equation. The value of $r$ corresponding to each value of $\rho$ is then found from one of the equations (I), most conveniently from the second as regards its absolute value. If $m+n$ is even, however, one of the other two equations must also be used to decide the sign.

\FifthArticle{19}
The indirect solution of equation II can always be carried out easily. Taking the following observations into account will help further.

1) The values of $\rho$ lie between $0$ and $180^{0}$. When the coefficients of the given equation are real, only half the number need be sought individually, namely those lying between $0$ and $90^{0}$.

2) In either case, first divide the relevant interval into the subdivisions formed by changes of sign in the values of the function of $\rho$ on the right-hand side of equation II. The transition values of $\rho$ can evidently occur only where one of the angles $m\rho+\varepsilon-\varphi,\ n\rho-\varepsilon,\ (m+n)\rho-\varphi$ is divisible by $180^{0}$, so that the function itself becomes either $0$ or infinite. Those subdivisions in which the value of the function becomes negative are then excluded of themselves from further investigation.

3) Unless approximate values of $\rho$ can already be obtained in other ways, the indirect search through the appropriate intervals will be made much easier by carrying out the first trials with shortened tables having few digits, as illustrated by the examples of Article 16. In some cases it may be convenient
\SourceLeaf{106}
to begin by using only logarithms of sines with three digits, written on a small sheet, say at intervals of one degree.

